\documentclass{article}
\usepackage[utf8]{inputenc}
\usepackage[spanish]{babel}
\usepackage{listings}
\usepackage{xcolor}
\usepackage{geometry}
\geometry{a4paper, margin=1in}

\% Configuración para código C++
\definecolor{codegreen}{rgb}{0,0.6,0}        
\definecolor{codegray}{rgb}{0.5,0.5,0.5}     
\definecolor{codepurple}{rgb}{0.58,0,0.82}   
\definecolor{backcolour}{rgb}{0.95,0.95,0.92}

\lstset{
    backgroundcolor=\color{backcolour},      
    commentstyle=\color{codegreen},
    keywordstyle=\color{magenta},
    numberstyle=\tiny\color{codegray},       
    stringstyle=\color{codepurple},
    basicstyle=\ttfamily\small,
    breakatwhitespace=false,
    breaklines=true,
    captionpos=b,
    keepspaces=true,
    numbers=left,
    numbersep=5pt,
    showspaces=false,
    showstringspaces=false,
    showtabs=false,
    tabsize=2,
    language=C++,
    inputencoding=utf8,
    extendedchars=true,
    literate={á}{{\'a}}1 {é}{{\'e}}1 {í}{{\'i}}1 {ó}{{\'o}}1 {ú}{{\'u}}1 {ñ}{{\~n}}1 {¿}{{\textquestiondown}}1 {¡}{{\textexclamdown}}1
}

\% Macros para las secciones de los apuntes
\newcommand{\quees}[1]{
    \section*{🤔 ¿QUÉ ES?}
    \hrulefill \\
    \#1
}

\newcommand{\dichodeotraforma}[1]{
    \section*{💬 DICHO DE OTRA FORMA}
    \hrulefill \\
    \#1
}

\newcommand{\diagrama}[1]{
    \section*{🎨 DIAGRAMA}
    \hrulefill \\
    \#1
}

\newcommand{\comofunciona}[1]{
    \section*{⚙️ ¿CÓMO FUNCIONA?}
    \hrulefill \\
    \#1
}

\newcommand{\complejidad}[1]{
    \section*{📱 COMPLEJIDAD (Big-O)}
    \hrulefill \\
    \#1
}

\newcommand{\entrevistas}[1]{
    \section*{🎯 EN ENTREVISTAS}
    \hrulefill \\
    \#1
}

\newcommand{\resumen}[1]{
    \section*{📋 RESUMEN RÁPIDO}
    \hrulefill \\
    \#1
}

\begin{document}

\title{Representación de Grafos}
\author{Aldo Zepeda Montoya}
\date{\today}
\maketitle

\subsection*{CATEGORÍA: 03\_Grafos}

\quees{
Un mapa de ciudades 🗺️ con carreteras entre ellas. 

En computación, las ciudades se llaman Nodos (o Vértices) y las 
carreteras se llaman Aristas (o Edges). Un grafo es simplemente un 
conjunto de nodos conectados entre sí. 

Puedes tener:
- Grafos Dirigidos: Carreteras de un solo sentido (A → B).
- Grafos No Dirigidos: Carreteras de doble sentido (A ↔ B).
- Grafos Ponderados: Carreteras con costo o distancia (A --5--> B).
}

\dichodeotraforma{
Un Grafo G = (V, E) es una estructura no lineal. Para trabajar con 
ellos en código, necesitamos guardarlos en memoria de alguna forma. 

Las dos formas más comunes son:
1. Adjacency Matrix (Matriz de Adyacencia): Una tabla de V x V.
2. Adjacency List (Lista de Adyacencia): Un array de listas.

La elección depende de si el grafo es "Denso" (muchas conexiones) o 
"Disperso" (pocas conexiones).
}

\diagrama{\begin{verbatim}
Grafo de ejemplo: A-B, A-C, B-C, C-D

Adjacency Matrix (4x4):      Adjacency List:
     A  B  C  D              A: [B, C]
  A [0, 1, 1, 0]             B: [A, C]
  B [1, 0, 1, 0]             C: [A, B, D]
  C [1, 1, 0, 1]             D: [C]
  D [0, 0, 1, 0]

- En la Matrix: 1 significa que hay conexión, 0 que no.
- En la List: Cada nodo tiene su propia lista de vecinos.
\end{verbatim}}

\begin{lstlisting}[language=C++]
1. Adjacency Matrix:
   - Pros: Muy rápido O(1) para saber si A y B están conectados.
   - Contras: Usa mucha memoria O(V^2), incluso si hay pocos edges.
   - Cuándo usar: Grafos pequeños o muy densos.

2. Adjacency List:
   - Pros: Ahorra mucha memoria O(V + E). Es la más usada.
   - Contras: Saber si A y B están conectados puede tardar O(V).
   - Cuándo usar: Casi siempre (la mayoría de los grafos reales son
     dispersos).

📝 PSEUDOCÓDIGO
──────────────────────────────────────────────────────────────

// Agregar arista en Lista de Adyacencia (no dirigido)
función agregarArista(u, v):
    adj[u].agregar(v)
    adj[v].agregar(u)

// Agregar arista en Matriz de Adyacencia
función agregarAristaMatrix(u, v):
    matrix[u][v] = 1
    matrix[v][u] = 1

💻 CÓDIGO C++
──────────────────────────────────────────────────────────────

#include <iostream>
#include <vector>
using namespace std;

int main() {
    int V = 4; // Número de vértices
    
    // 1. Representación con Lista de Adyacencia (Recomendado)
    vector<vector<int>> adj(V);

    // Agregar A-B (0-1), A-C (0-2), B-C (1-2), C-D (2-3)
    adj[0].push_back(1); adj[1].push_back(0);
    adj[0].push_back(2); adj[2].push_back(0);
    adj[1].push_back(2); adj[2].push_back(1);
    adj[2].push_back(3); adj[3].push_back(2);

    cout << "Lista de Adyacencia de C (nodo 2): ";
    for(int vecino : adj[2]) {
        cout << vecino << " ";
    }
    cout << endl;

    // 2. Representación con Matriz de Adyacencia
    vector<vector<int>> matrix(V, vector<int>(V, 0));
    matrix[0][1] = 1; matrix[1][0] = 1;
    // ... completar resto

    return 0;
}

⏱️ COMPLEJIDAD (Comparativa)
──────────────────────────────────────────────────────────────

  Operación        │ Matrix   │ List
  ─────────────────┼──────────┼──────────────
  Espacio          │ O(V^2)    │ O(V + E)
  Agregar Edge     │ O(1)     │ O(1)
  Verificar Edge   │ O(1)     │ O(Grado del nodo)
  Iterar Vecinos   │ O(V)     │ O(Grado del nodo)

* V = Vértices, E = Edges (Aristas).
\end{lstlisting}

\entrevistas{
✅ Default: En el 90\% de las entrevistas, usa Adjacency List. Es más
   eficiente en memoria y espacio.

💡 Representación con Map: A veces los nodos no son números (0, 1, 2)
   sino nombres ("CDMX", "NY"). En ese caso, usa:
   `unordered\_map<string, vector<string>> adj;`

⚠️ Edge List: Existe una tercera forma llamada "Edge List" que es 
   simplemente un array de pares `[(0,1), (0,2), ...]`. Útil para 
   algoritmos como Kruskal.

🔑 Grafo Implícito: A veces el grafo no te lo dan, sino que es una
   matriz (grid) donde puedes moverte arriba, abajo, izq, der. Ahí el
   "grafo" son las celdas y sus vecinos.
}

\resumen{
• Nodos = Ciudades, Aristas = Caminos.
• Matrix: Rápida pero traga mucha memoria O(V^2).
• List: Eficiente en memoria, la favorita de los devs O(V + E).
• Dirigido (flecha) vs No Dirigido (línea).
• Ponderado (tiene peso/costo) vs No Ponderado.
• Grado de un nodo = cuántos vecinos tiene.
════════════════════════════════════════════════════════════════
}

\end{document}